\documentclass[11pt]{article}

%------------------------- fonts & encoding -----------------------------------
\usepackage[T1]{fontenc}
\usepackage[utf8]{inputenc}
\usepackage{amsmath,amssymb,amsthm}
\usepackage{mathtools}
\usepackage[sc]{mathpazo}
\linespread{1.06}
\usepackage{microtype}

%------------------------- layout ---------------------------------------------
\usepackage[a4paper,margin=1.05in,marginparwidth=1.0in]{geometry}
\usepackage{booktabs}
\usepackage{enumitem}
\setlist{itemsep=2pt,topsep=3pt}
\usepackage{xcolor}
\usepackage{graphicx}
\usepackage{tikz}
\usepackage{pgfplots}
\pgfplotsset{compat=1.17}
\usepackage[most,breakable]{tcolorbox}

%------------------------- palette --------------------------------------------
\definecolor{crimson}{HTML}{A51C30}
\definecolor{ink}{HTML}{1A1A1A}
\definecolor{steel}{HTML}{2C5F8A}
\definecolor{amber}{HTML}{9A6A00}
\definecolor{teal}{HTML}{1F6F6B}
\definecolor{paperlt}{HTML}{F7F4EF}

%------------------------- section styling ------------------------------------
\usepackage{titlesec}
\titleformat{\section}{\Large\bfseries\color{steel}}{\thesection}{0.7em}{}
\titleformat{\subsection}{\large\bfseries\color{ink}}{\thesubsection}{0.6em}{}
\titlespacing*{\section}{0pt}{1.6em}{0.6em}

%------------------------- hyperref -------------------------------------------
\usepackage[colorlinks=true,linkcolor=steel,citecolor=steel,urlcolor=crimson]{hyperref}

%------------------------- theorem environments -------------------------------
\theoremstyle{definition}
\newtheorem{definition}{Definition}[section]
\newtheorem{example}{Example}[section]
\newtheorem{exercise}{Exercise}[section]
\theoremstyle{plain}
\newtheorem{proposition}{Proposition}[section]
\theoremstyle{remark}
\newtheorem*{remark}{Remark}

%------------------------- callout boxes --------------------------------------
\newtcolorbox{intuition}{
  enhanced, breakable, colback=crimson!5, colframe=crimson, boxrule=0.5pt,
  left=8pt,right=8pt,top=6pt,bottom=6pt, arc=2pt,
  colbacktitle=crimson, coltitle=paperlt, fonttitle=\bfseries, title={Intuition}}

\newtcolorbox{pitfall}{
  enhanced, breakable, colback=amber!7, colframe=amber, boxrule=0.5pt,
  left=8pt,right=8pt,top=6pt,bottom=6pt, arc=2pt,
  colbacktitle=amber, coltitle=paperlt, fonttitle=\bfseries, title={Pitfall}}

\newtcolorbox{keyresult}{
  enhanced, breakable, colback=steel!6, colframe=steel, boxrule=0.5pt,
  left=8pt,right=8pt,top=6pt,bottom=6pt, arc=2pt,
  colbacktitle=steel, coltitle=paperlt, fonttitle=\bfseries, title={Takeaway}}

\newtcolorbox{observed}{
  enhanced, breakable, colback=teal!6, colframe=teal, boxrule=0.5pt,
  left=8pt,right=8pt,top=6pt,bottom=6pt, arc=2pt,
  colbacktitle=teal, coltitle=paperlt, fonttitle=\bfseries, title={Measured}}

%------------------------- math operators -------------------------------------
\DeclareMathOperator{\Var}{Var}
\DeclareMathOperator{\Cov}{Cov}
\DeclareMathOperator{\Corr}{corr}
\DeclareMathOperator{\diag}{diag}
\DeclareMathOperator{\std}{std}
\DeclareMathOperator*{\avg}{\mathbb{E}}
\newcommand{\T}{^{\!\top}}
\newcommand{\R}{\mathbb{R}}
\newcommand{\half}{\tfrac12}
\newcommand{\wvec}{\mathbf{w}}
\newcommand{\Xmat}{\mathbf{X}}
\newcommand{\Fmat}{\mathbf{F}}
\newcommand{\Vmat}{\mathbf{V}}
\newcommand{\Umat}{\mathbf{U}}
\newcommand{\Rmat}{\mathbf{R}}
\newcommand{\Lam}{\boldsymbol{\Lambda}}
\newcommand{\Del}{\boldsymbol{\Delta}}
\newcommand{\fvec}{\mathbf{f}}

%------------------------- title ----------------------------------------------
\title{\vspace{-2.2em}
  {\color{steel}\rule{\textwidth}{1.3pt}}\\[0.7em]
  \Huge\bfseries Liquidity\\[0.35em]
  \large\mdseries The trading-activity factor --- share turnover at three horizons, on a log scale\\[0.4em]
  {\color{steel}\rule{\textwidth}{0.7pt}}}
\author{}
\date{}

%------------------------- watermark footer -----------------------------------
\usepackage{fancyhdr}
\pagestyle{fancy}
\fancyhf{}
\renewcommand{\headrulewidth}{0pt}
\renewcommand{\footrulewidth}{0pt}
\fancyfoot[L]{\footnotesize\color{black!45}MSCI Barra USE4 Imitation}
\fancyfoot[R]{\footnotesize\color{black!45}\thepage}

\begin{document}
\maketitle
\thispagestyle{empty}

\begin{abstract}
Liquidity is the trading-activity factor: it measures how actively a stock changes hands, so that heavily-traded names load high and sleepy ones load low. USE4 measures it as \emph{share turnover} -- the fraction of a firm's shares that trade in a given window -- read at three horizons (one month, one quarter, one year) and combined into a single descriptor. Two construction choices dominate the factor. First, turnover is divided by shares outstanding rather than left as a dollar figure, so the measure is the fraction of the company that trades, comparable across a giant and a micro-cap alike. Second, and most importantly, the descriptor is built in \emph{logs}: raw turnover is wildly right-skewed, spanning orders of magnitude from a sleepy utility to a frenzied meme stock, and the logarithm compresses that into a nearly symmetric, well-behaved quantity. A side effect of the log is that the raw descriptor is \emph{negative} for almost every firm -- turnover is a fraction below one, and the log of a fraction is negative -- which is a feature of the scale, not a sign of illiquidity. The result is the calmest of the style factors: after logging, the outlier trim barely bites and the standardised exposure is close to Gaussian.
\end{abstract}

\begin{tcolorbox}[colback=paperlt,colframe=ink!30,boxrule=0.4pt,arc=2pt,left=10pt,right=10pt,top=8pt,bottom=8pt]
\textbf{How to read these notes.} We move from \emph{why} Liquidity is a factor to \emph{how} USE4 specifies it, then to how we actually compute turnover at three horizons, with a deliberate stop at the step that makes the factor work -- the logarithm (Section~\ref{sec:log}) -- and at how to read a descriptor that is negative for nearly everyone (Section~\ref{sec:reading}). Read those two slowly; they carry the weight of the set. The numbers labelled as measured come from a single run of our personal build.

\smallskip
Four coloured boxes recur: \textcolor{crimson}{\textbf{Intuition}} gives the mental picture, \textcolor{amber}{\textbf{Pitfall}} flags what breaks without the fix, \textcolor{steel}{\textbf{Takeaway}} states a load-bearing result, and \textcolor{teal}{\textbf{Measured}} reports real numbers from the build run.

\smallskip
\textbf{Reference material.}
\begin{itemize}[leftmargin=1.4em,topsep=2pt]
  \item External: Menchero, Orr \& Wang (2011), USE4 Empirical Notes (Appendix A) and Methodology Notes (Sec. 2.2--2.3).
\end{itemize}
\end{tcolorbox}

\tableofcontents
\bigskip

\section{What Liquidity is, and why it is a factor}\label{sec:what}

\paragraph{Learning objectives.}
\begin{itemize}[leftmargin=1.4em,topsep=2pt]
  \item State Liquidity as trading activity, measured by share turnover.
  \item Explain why turnover is divided by shares outstanding rather than left in dollars.
  \item Say in one sentence what the Liquidity exposure is, before any standardisation.
\end{itemize}

Liquidity is about how easily a stock trades. Some names change hands constantly -- a large fraction of their shares turning over every month -- while others sit quietly, with only a sliver trading. That difference matters for risk in two ways. A thinly-traded stock is costly and slow to exit, so investors demand a premium to hold it; and trading activity tends to move together across the market, surging in booms and drying up in stress, so a liquidity shock hits the quiet names as a group. A risk model needs to see that shared exposure directly.

The natural yardstick is \emph{share turnover}: the number of shares traded over some window, divided by the shares outstanding. Dividing by shares outstanding is the crucial move. Raw dollar volume would mostly just measure size -- big companies trade more dollars simply because they are big -- whereas turnover asks what \emph{fraction of the company} changed hands, which is comparable across a trillion-dollar giant and a small-cap alike. A turnover of $0.20$ means a fifth of the firm's shares traded in the window, whether the firm is huge or tiny.

\begin{intuition}
The whole factor answers one question: \emph{what fraction of this company trades, and how often?} A high-turnover name is liquid and closely watched; a low-turnover name is quiet and harder to move in size. Hold that picture -- a \emph{fraction of the float}, not a dollar amount -- and the horizons and the logarithm below are just refinements of it.
\end{intuition}

\begin{keyresult}
The Liquidity exposure is a standardised blend of log share turnover at three horizons. Everything that follows is about which horizons, and why the logarithm is not optional.
\end{keyresult}

\section*{Notation}
\addcontentsline{toc}{section}{Notation}

\begin{center}
\begin{tabular}{@{}ll@{}}
\toprule
Symbol & Meaning \\
\midrule
$V_{i,d}$ & shares of stock $i$ traded on day $d$ (daily volume) \\
$S_{i,d}$ & shares outstanding of stock $i$ on day $d$ \\
$V_{i,d}/S_{i,d}$ & daily share turnover (fraction of the float traded) \\
$\text{STOM}_i$ & log one-month share turnover, $\ln\!\big(\sum_{d\in\text{month}} V_{i,d}/S_{i,d}\big)$ \\
$\text{STOQ}_i$ & log average monthly turnover over the past quarter (3 months) \\
$\text{STOA}_i$ & log average monthly turnover over the past year (12 months) \\
$c_i$ & raw LIQ composite, $0.35\,\text{STOM}_i+0.35\,\text{STOQ}_i+0.30\,\text{STOA}_i$ \\
$x_i$ & standardised LIQ exposure for stock $i$ \\
$w_i$ & cap weight of stock $i$ in the ESTU \\
$\avg_{\text{cw}}, \std_{\text{ew}}$ & cap-weighted mean, equal-weighted standard deviation over the ESTU \\
$Q_1,\dots,Q_5$ & quintiles of raw $c_i$ (all negative), from least to most traded \\
ESTU & estimation universe (MSCI USA IMI proxy) \\
\bottomrule
\end{tabular}
\end{center}

\section{How USE4 defines it}\label{sec:def}

\paragraph{Learning objectives.}
\begin{itemize}[leftmargin=1.4em,topsep=2pt]
  \item Name the three turnover horizons USE4 blends and their weights.
  \item State that each is a \emph{log} turnover, and the standardisation target.
\end{itemize}

USE4 builds Liquidity from three share-turnover descriptors (Menchero, Orr \& Wang 2011, Empirical Notes, App. A): STOM, the log of one month's share turnover; STOQ, the log of average monthly turnover over the past quarter; and STOA, the log of average monthly turnover over the past year. The three are near co-equal, weighted $0.35$, $0.35$, and $0.30$ -- the recent month and the recent quarter share top billing, with the full-year average close behind. All three are measured on the same scale (monthly turnover) and differ only in how far back they look.

Each descriptor is a \emph{logarithm} of turnover, not turnover itself -- the reason is the whole subject of Section~\ref{sec:log}. The weighted composite is then trimmed at three standard deviations (Methodology Notes, Sec. 2.2) and standardised over the estimation universe to cap-weighted mean $0$ and equal-weighted standard deviation $1$ (Sec. 2.3). Turnover is computed from daily data aligned point-in-time, so any signal date uses only trading a market participant could already have observed.

\section{How we compute it: turnover at three horizons}\label{sec:compute}

\paragraph{Learning objectives.}
\begin{itemize}[leftmargin=1.4em,topsep=2pt]
  \item Reconstruct STOM, STOQ, and STOA from daily volume and shares outstanding.
  \item Explain why three horizons are blended rather than one.
  \item Describe how a newer stock with a short history still earns a score.
\end{itemize}

The raw ingredient is daily share turnover: each day's trading volume divided by that day's shares outstanding, $V_{i,d}/S_{i,d}$, the fraction of the float that changed hands. Daily volume and a daily share count both come from the daily price-and-volume record (Sharadar SEP), aligned point-in-time so no future trading leaks into a signal date. Summing the daily turnovers across roughly a month of trading days gives the month's total turnover; \textbf{STOM} is its logarithm. \textbf{STOQ} and \textbf{STOA} take the same monthly turnover and average it over the past three and twelve months respectively, then take the log -- smoother, slower-moving views of the same quantity.

Why three horizons rather than one? Because trading activity is both spiky and persistent. A single month can be distorted by a one-off event -- an index addition, an earnings surprise, a news frenzy -- that sends turnover soaring for a few weeks. STOM captures that recent burst; STOA captures the firm's sustained baseline liquidity over a full year; STOQ sits between. Blending the three balances responsiveness against stability, so the factor reacts to genuine changes in liquidity without overreacting to a single noisy month.

A few practical choices keep the measure honest. A day with zero volume contributes zero turnover, not a gap -- a stock that did not trade is genuinely illiquid that day, and the zero says so; a month with no trading at all yields no valid turnover for that window rather than a spurious value. A newer stock without a full year of history is not discarded: as long as it has a valid recent month it earns a score, and the blend simply renormalises its weights over the horizons that do exist, so a young firm with a few months of data contributes the horizons it can support rather than dropping out entirely.

\begin{pitfall}
Get the denominator and the units right. Turnover must divide by \emph{shares outstanding}, not be left as dollar volume -- otherwise the factor silently re-measures size, since big firms trade more dollars by default. And the daily volume-to-shares ratio must be on a consistent unit basis; a units mismatch between a volume field and a shares field can inflate turnover by a constant factor across the board, a bug that hides because every name moves together and the cross-sectional ranking still looks plausible.
\end{pitfall}

\section{Why take logs: taming a skewed quantity}\label{sec:log}

\paragraph{Learning objectives.}
\begin{itemize}[leftmargin=1.4em,topsep=2pt]
  \item Explain why raw turnover is unusable as a linear descriptor.
  \item Describe what the logarithm does to the distribution, and why that makes Liquidity well-behaved.
\end{itemize}

This is the step that makes Liquidity work. Raw turnover is one of the most skewed quantities in any equity dataset: a sleepy stock might trade a few percent of its float a month, while a stock in a speculative frenzy trades several times its entire float in the same span. Those differ not by a little but by orders of magnitude. On a linear scale the frenzied names dominate any average, the cross-section is all right tail, and a three-standard-deviation rule would either clip a huge fraction of names or let a handful of outliers set the scale for everyone.

The logarithm fixes this at a stroke. Turnover is a fundamentally \emph{multiplicative} quantity -- one stock trades ten times another -- and the log turns multiplicative spread into additive spread, pulling the long right tail back in and stretching the compressed low end out. The result is a nearly symmetric, roughly bell-shaped distribution (Figure~\ref{fig:liq_dist}). That is why Liquidity is the best-behaved of the style factors: once the descriptor is in logs, the outlier trim barely has anything to clip, and the standardised exposure comes out close to Gaussian.

\begin{figure}[htbp]
\centering
\includegraphics[width=\linewidth]{figures/liq_dist.png}
\caption{Why Liquidity is built in logs. \emph{Left:} raw blended monthly turnover, pooled over in-ESTU stock-dates, is violently right-skewed -- a dense pile of low-turnover names and a tail of hyperactive ones that runs off the axis. \emph{Right:} the same quantity in logs -- the LIQ descriptor -- is nearly symmetric and bell-shaped, centred near $-1.8$. The log turns an unusable multiplicative spread into a clean additive one; it is also why the descriptor is negative (the log of a fraction below one).}
\label{fig:liq_dist}
\end{figure}

\begin{intuition}
Logs are the right language for anything that varies multiplicatively -- incomes, city sizes, trading volumes. Asking ``how many times more does A trade than B?'' is the natural question for turnover, and the log answers it on an additive ruler where equal distances mean equal \emph{ratios}. A name one unit of log-turnover above another trades about $e \approx 2.7$ times as much, anywhere along the scale.
\end{intuition}

\section{Reading the cross-section: negative by construction}\label{sec:reading}

\paragraph{Learning objectives.}
\begin{itemize}[leftmargin=1.4em,topsep=2pt]
  \item Explain why the raw descriptor is negative for nearly every firm, and why that is harmless.
  \item Contrast the in-universe distribution with the all-stocks distribution.
\end{itemize}

A reader meeting the raw descriptor for the first time is often startled that it is negative almost everywhere -- the median is around $-1.8$. This is pure arithmetic, not a finding. Monthly turnover is a fraction below one for the overwhelming majority of stocks (few firms trade more than their whole float in a month), and the logarithm of a number below one is negative; the rare exceptions are a handful of frenzied names that trade more than their entire float in a month, which carry a small positive descriptor. So the negative values carry no message about ``illiquidity'' in absolute terms; what matters is a stock's position \emph{relative to the market}, which is exactly what the later standardisation supplies. After standardising, the exposure is centred at zero and the sign of the standardised score -- above or below the dollar-weighted market -- is what you read.

There is also a striking gap between the investable universe and the market as a whole. Within the estimation universe -- the liquid, investable core -- the standardised exposure is tightly behaved, with almost every name inside the usual $[\pm 3]$ band. Across \emph{all} stocks, far fewer are: the full universe contains an enormous tail of barely-traded micro-caps whose turnover sits far below the investable core, well out in the left tail of the standardised scale. Because the standardisation is calibrated on the estimation universe, those out-of-universe names land far outside the band. The contrast is a useful reminder that the estimation universe is the liquid heart of a much less liquid market.

\begin{observed}
On the build run the raw liquidity descriptor (log turnover) has an ESTU median of $\approx -1.8$ and is nearly symmetric -- skew $\approx +0.1$ and close to Gaussian (excess kurtosis only $\approx 0.6$, a dramatic reduction from the raw turnover's tail), a direct result of the log transform. The three-sigma trim clips only $\approx 0.95\%$ of ESTU rows, comfortably inside the usual $0.5$--$2\%$ band and the gentlest trim of any style factor. After standardisation the cap-weighted mean is $0$ to machine precision, the equal-weighted standard deviation is $1.0$, and the standardised median is essentially zero (the distribution is symmetric). About $99\%$ of ESTU rows fall inside $[\pm 3]$, but only $\approx 84\%$ of \emph{all} scored stocks do -- the out-of-universe illiquid tail sits far below the investable core. Quintiles of raw turnover are strictly monotone (and all negative, on the log scale), from $Q_1 \approx -3.0$ up to $Q_5 \approx -0.6$, a spread of $\approx 2.35$. Month-over-month rank stability is high, Spearman $\rho \approx 0.97$. The deliverable is $\approx 1.75$M scored rows over $330$ dates and $\approx 16{,}580$ tickers.
\end{observed}

\section{Trim and standardisation}\label{sec:std}

\paragraph{Learning objectives.}
\begin{itemize}[leftmargin=1.4em,topsep=2pt]
  \item State the standardisation target and explain why the trim does so little work here.
\end{itemize}

The raw composite $c_i$ is trimmed at three standard deviations and standardised over the ESTU to cap-weighted mean $0$ and equal-weighted standard deviation $1$:
\[
x_i \;=\; \frac{c_i - \avg_{\text{cw}}(c)}{\std_{\text{ew}}(c)},
\qquad
\avg_{\text{cw}}(c) = \frac{\sum_i w_i\, c_i}{\sum_i w_i}.
\]
The cap-weighted centring measures a stock's trading activity against the dollar-weighted market; the equal-weighted scaling fixes the spread. Because the log has already pulled in the tails, the trim removes only a sliver -- under one percent of ESTU rows -- where for a raw-turnover descriptor it would have had to clip aggressively. The log does the heavy lifting; the trim merely tidies the edges.

\begin{exercise}
A stock trades, on average, $0.05$ of its float per month (5\% monthly turnover). Compute its STOM (use $\ln 0.05 \approx -3.0$). Is the result negative, and does that mean the stock is ``illiquid'' in any absolute sense, or only relative to the market once standardised?
\end{exercise}

\begin{exercise}
Two stocks have raw log-turnover descriptors that differ by exactly $1.0$ before standardisation. By roughly what factor do their turnovers differ? (Hint: the descriptor is a natural log.)
\end{exercise}

\begin{exercise}
Liquidity clips only $\approx 0.95\%$ of names at three standard deviations, while a balance-sheet factor like leverage clips $\approx 2.7\%$. Explain in one or two sentences why a log-transformed turnover descriptor needs so little trimming, referring to the shape of the distribution in Figure~\ref{fig:liq_dist}.
\end{exercise}

\section{Summary, and one caveat}\label{sec:summary}

\paragraph{Learning objectives.}
\begin{itemize}[leftmargin=1.4em,topsep=2pt]
  \item Recite the Liquidity build from daily volume to standardised exposure.
  \item State the central limitation and what would address it.
\end{itemize}

The build, end to end: form daily share turnover as volume over shares outstanding; sum it across a month and average that monthly turnover over a quarter and a year; take logs to get STOM, STOQ, and STOA; blend them $0.35/0.35/0.30$, renormalising the weights for a young stock that lacks the longer histories; trim the composite at three standard deviations -- which barely bites, because the log already tamed the skew -- and standardise to cap-weighted mean $0$, equal-weighted standard deviation $1$. The output is a near-symmetric, well-behaved trading-activity factor, negative on the raw log scale but centred at zero once standardised.

\begin{keyresult}
Liquidity is share turnover, read at three horizons and measured in logs. The two computations that matter most are dividing by shares outstanding -- so the factor measures activity, not size -- and the logarithm, which turns a wildly skewed quantity into the calmest descriptor in the model.
\end{keyresult}

\begin{pitfall}
One honest limitation matters most. Turnover measures trading \emph{activity}, which is not the same as trading \emph{cost}: a stock in a speculative frenzy shows enormous turnover yet can be treacherous to exit in size, while a steadily-held quality name may trade modestly and yet be cheap to transact. High turnover therefore does not guarantee low transaction cost -- the factor captures how much a name trades, not how gently a large order would move its price. A smaller, second limitation is that turnover here divides by total shares outstanding rather than free float, which understates turnover for closely-held firms whose tradable float is much smaller than their share count. The clear levers are a free-float share count for the denominator, and a cost-based liquidity measure (such as a bid-ask or price-impact statistic) to complement the activity-based one.
\end{pitfall}

\begin{remark}
None of this changes the factor's place in the model: Liquidity remains the trading-activity dimension, sitting alongside size, value, and the rest. The caveat is about \emph{what} turnover can and cannot stand in for, not about whether the activity signal is real.
\end{remark}

\begin{center}\small\itshape
References: Menchero, Orr \& Wang (2011), \textit{The Barra US Equity Model (USE4)} --- Empirical Notes (Appendix A) and Methodology Notes (Sec. 2.2--2.3). ESTU treated as an MSCI USA IMI proxy; measured quantities are from a personal build run.
\end{center}

\end{document}
